\section{Embodied-Motivational Architecture for Autonomous Light Seeking}
\label{sec:embodied_motivational_light_seeker}

\subsection{Scope}
This section operationalizes two foundational branches from \texttt{cartheur/ideal} for a Bluetooth-connected mobile robot tasked with seeking the brightest light source:
\begin{enumerate}
\item \textbf{01-embodied-paradigm}: cognition is organized around \textit{experiment $\rightarrow$ result} interactions instead of passive world-state observation.
\item \textbf{02-interactional-motivation}: interaction outcomes are assigned valence, shaping future experiment selection.
\end{enumerate}

\subsection{Formalization}
Let $\mathcal{E}$ be a finite set of experiments (e.g., \texttt{scan-left}, \texttt{scan-right}) and $\mathcal{R}$ be a finite set of results (e.g., \texttt{dark}, \texttt{bright}).
At cycle $t$, the agent executes:
\begin{equation}
e_t \in \mathcal{E}, \quad r_t = f(e_t)
\end{equation}
where $f$ is the embodied environment coupling.

The agent stores primitive interactions:
\begin{equation}
I_t = \langle e_t, r_t \rangle
\end{equation}
and predicts anticipated result $\hat{r}(e)$ from memory.

Motivational valence is modeled as:
\begin{equation}
v: \mathcal{R} \rightarrow \mathbb{R}
\end{equation}
with $v(\texttt{bright}) > 0$ and $v(\texttt{dark}) < 0$ for this mission.

The action-selection rule is:
\begin{equation}
e_{t+1} = \arg\max_{e \in \mathcal{E}} v(\hat{r}(e))
\end{equation}
when predictions are available, with a boredom-driven exploration fallback.

\subsection{Implementation Mapping}
The \texttt{haden} implementation introduces:
\begin{itemize}
\item \texttt{Haden.Library/Algorithm/IdealEmbodiedLightSeeker.cs}
\item \texttt{Haden.Library/Algorithm/IdealEmbodiedDecisionTreeAdapter.cs}
\item \texttt{Haden.Tests/IdealBranch12LightSeekerTests.cs}
\end{itemize}

The model tracks:
\begin{enumerate}
\item primitive interaction memory;
\item anticipation correctness;
\item mood transitions (\texttt{Frustrated}, \texttt{SelfSatisfied}, \texttt{Bored}, \texttt{Pleased}, \texttt{Pained});
\item boredom-based exploration and valence-based exploitation.
\end{enumerate}

\subsection{Mission Interpretation}
For the light-seeker mission, each scan-direction command is an experiment and each normalized light reading class is a result.
The policy converges toward experiments anticipating positive-valence outcomes, operationalized as bright-source alignment.
This framing supports both simulation and physical execution over Bluetooth serial transport without changing the core cognition loop.

\subsection{Testable Claims}
\begin{enumerate}
\item The branch-01 embodied loop exhibits reproducible boredom-triggered experiment switching.
\item The branch-02 motivational extension converges toward bright-result experiments under stable valence assignment.
\end{enumerate}

These claims are encoded as executable regression tests in \texttt{IdealBranch12LightSeekerTests}.

\subsection{Academic Value}
This formulation contributes three research-relevant properties:
\begin{enumerate}
\item \textbf{Construct validity}: branch-01 and branch-02 theoretical claims are mapped to explicit computational observables (interaction tuples, anticipation outcomes, mood state, and valence-driven experiment choice).
\item \textbf{Reproducibility}: deterministic tests provide executable evidence for claimed behavioral regimes, reducing ambiguity between conceptual prose and implementation.
\item \textbf{Externalization path}: the same embodied loop can run in deterministic simulation and on Bluetooth-connected hardware, enabling controlled ablation in simulation followed by ecological evaluation on the physical robot.
\end{enumerate}

For paper reporting, this supports a methods section that cleanly separates:
\begin{enumerate}
\item theory (embodied interaction and motivation),
\item implementation (state variables and update equations), and
\item evaluation (trace-level behavioral criteria and mission outcomes).
\end{enumerate}

\subsection{Live Robot Seek-Max-Light Test Protocol}
\label{sec:live_seek_max_light_protocol}

\textbf{Objective.}
Validate that the migrated Linux autonomy routine preserves the latest legacy behavior while physically orienting the NXT toward increasing light intensity.

\textbf{System under test.}
\begin{itemize}
\item decision engine: \texttt{Haden.RobotBehavior/LegacyAutonomousLightSeekEngine.cs};
\item live runner: \texttt{Haden.HardwareSmoke} in \texttt{--seek-max-light} mode;
\item transport: Bluetooth RFCOMM (default \texttt{/dev/rfcomm0}).
\end{itemize}

\textbf{Hardware and environment.}
\begin{enumerate}
\item LEGO NXT brick paired with Linux host;
\item light sensor connected to configured input port (default Port 3);
\item seek motor connected to configured output port (default Port A);
\item single fixed light source in controlled ambient illumination.
\end{enumerate}

\textbf{Procedure.}
For each iteration $t \in \{1,\dots,N\}$:
\begin{enumerate}
\item read sensor value $s_t$;
\item advance legacy autonomous state machine (\texttt{Peek/Compare}-style update);
\item apply turn command ($\texttt{Left}$ or $\texttt{Right}$) with fixed power and angular step;
\item wait a fixed settle delay and repeat.
\end{enumerate}
After $N$ iterations, brake motor and disconnect.

\textbf{Trace variables.}
Per iteration, log: sensor reading, chosen turn direction, current value, peak value, and reward counter.

\textbf{Primary evaluation criteria.}
\begin{enumerate}
\item peak value increases monotonically or stabilizes near a local maximum;
\item direction reversal occurs on negative gradient events;
\item final steady-state reading exceeds initial reading under identical lighting.
\end{enumerate}

\textbf{Reproducibility controls.}
Fix source position, ambient light, motor power, turn granularity, and iteration count.
Repeat across multiple initial headings; report mean and variance of final peak and convergence steps.

\textbf{Reference execution command.}
\begin{verbatim}
HADEN_LIGHT_SENSOR_PORT=3 \
HADEN_LIGHT_SEEK_MOTOR_PORT=0 \
HADEN_SEEK_ITERATIONS=30 \
HADEN_SEEK_POWER=25 \
HADEN_SEEK_DEGREES=30 \
dotnet run --project Haden.HardwareSmoke/Haden.HardwareSmoke.csproj \
  -- --seek-max-light /dev/rfcomm0
\end{verbatim}
